%!TEX program = uplatex
\RequirePackage{plautopatch}

\documentclass[uplatex,dvipdfmx]{jlreq}

% 余白調整（1ページ収め）
\usepackage[top=18mm,bottom=20mm,left=20mm,right=20mm]{geometry}

\usepackage{amsmath,amssymb}
\usepackage{url}

\pagestyle{empty}

% 行間少し詰める
\renewcommand{\baselinestretch}{0.96}

%--------------------------------
% 講演マクロ
%--------------------------------

\newcommand{\talk}[3]{%
\noindent
#1\quad #2\par
\hspace*{2em}#3\par\smallskip
}

%--------------------------------
% タイトル情報
%--------------------------------

\newcommand{\EwMnumber}{第22回}
\newcommand{\EwMtitle}{「離散」の世界\\[2mm]\large 可換から非可換へ}
\newcommand{\EwMdate}{2002年2月1日（金）14:30～17:40 ～ 2月2日（土）10:30～16:50}
\newcommand{\EwMplace}{東京都文京区春日1-13-27 中央大学理工学部5号館 5533号室}

%--------------------------------
% タイトル表示
%--------------------------------

\newcommand{\EwMtitleblock}{

\vspace*{-5mm}

\begin{center}

{\Large ENCOUNTER with MATHEMATICS}

\vspace{2mm}

{\large 数学との遭遇}

\vspace{4mm}

{\Large \bfseries \EwMnumber\ \EwMtitle}

\vspace{4mm}

\EwMdate

\vspace{2mm}

於：\EwMplace

\end{center}

\vspace{2mm}

}

\begin{document}

\EwMtitleblock

%--------------------------------
% プログラム
%--------------------------------

\section*{プログラム}

\subsection*{2月1日（金）}

\talk
{14:30～15:30}
{たかがアーベル、されどアーベル I}
{砂田 利一 氏（東北大・理）}

\talk
{16:10～17:40}
{結晶格子上のランダム・ウォークと結晶格子の幾何}
{小谷 元子 氏（東北大・理）}

\subsection*{2月2日（土）}

\talk
{10:30～12:00}
{幾何学的群論の話}
{藤原 耕二 氏（東北大・理）}

\talk
{14:00～15:30}
{ユニタリ表現から見た非可換性\\[1mm]Kazhdan の性質 (T)}
{井関 裕泰 氏（東北大・理）}

\talk
{15:50～16:50}
{たかがアーベル、されどアーベル II}
{砂田 利一 氏（東北大・理）}

\bigskip

17:00～　懇親会（ワイン・パーティー）

%--------------------------------
% 案内文
%--------------------------------

\bigskip

別紙の趣旨に沿った集会の第22回を以上のような予定で開催いたします。  
非専門家向けに入門的な講演をお願い致しました。  
多く方々のご参加をお待ちしております。  
講演者による講演内容へのご案内を別紙に添付いたしますのでご覧下さい。

%--------------------------------
% 連絡先
%--------------------------------

\section*{連絡先}

112-8551 東京都文京区春日1-13-27 中央大学理工学部数学教室 

tel : 03-3817-1745  

ENCOUNTER with MATHEMATICS : e-mail : encmath@math.chuo-u.ac.jp  
homepage : \url{http://www.math.chuo-u.ac.jp/ENCwMATH}  

三松 佳彦 : yoshi@math.chuo-u.ac.jp / 高倉 樹 : takakura@math.chuo-u.ac.jp

\end{document}
